\documentclass[10pt]{article}
\usepackage[T1]{fontenc}
\usepackage[utf8]{inputenc}
\usepackage[french]{babel}
\usepackage{amsmath,amssymb,mathrsfs,mathpazo}
\usepackage{geometry}
\geometry{a4paper,margin=2cm}
\usepackage{graphicx,array,tabularx,tabularray,colortbl,multirow,multicol,booktabs,enumitem,pifont,ulem,fancybox,diagbox,cancel}
\usepackage{pstricks,pstricks-add,pst-plot,pst-tree,pst-node,pst-text,pst-func,pst-math,pst-bspline,pst-3dplot,pst-eucl,pst-solides3d}
\let\pscalculate\undefined
\usepackage{pst-bezier}
\usepackage{tikz,pgfplots,tkz-tab}
\usetikzlibrary{arrows,arrows.meta,patterns,calc,positioning,trees,decorations.markings,shapes.geometric,shapes.misc,shapes.arrows,intersections,decorations.pathmorphing,decorations.pathreplacing,backgrounds,shadows,babel}
\pgfplotsset{compat=1.18}
\usepgfplotslibrary{fillbetween}
\usepackage[np]{numprint}
\usepackage{listings,makecell,pas-tableur,twoopt}
\let\barre\undefined
\usepackage[pstricks]{tablvar}
\definecolor{darkgreen}{RGB}{0,100,0}
\providecommand{\red}{\color{red}}
\providecommand{\blue}{\color{blue}}
\providecommand{\green}{\color{green}}
\providecommand{\black}{\color{black}}
\providecommand{\white}{\color{white}}
\providecommand{\gray}{\color{gray}}
\providecommand{\cyan}{\color{cyan}}
\providecommand{\magenta}{\color{magenta}}
\providecommand{\yellow}{\color{yellow}}
\providecommand{\violet}{\color{violet}}
\providecommand{\orange}{\color{orange}}
\providecommand{\euro}{\texteuro}
\providecommand{\EUR}{\texteuro}
\providecommand{\ieme}{\textsuperscript{e}}
\providecommand{\degres}{\ensuremath{^\circ}}
\providecommand{\degree}{\ensuremath{^\circ}}
\providecommand{\vv}[1]{\overrightarrow{#1}}
\providecommand{\Coord}[1]{#1}
\providecommand{\decosix}{}
\providecommand{\starredbullet}{\ensuremath{\bullet}}
\pagestyle{empty}
\setlength{\parindent}{0pt}
\usepackage{adjustbox}\DeclareUnicodeCharacter{2236}{:}\DeclareUnicodeCharacter{2190}{\ensuremath{\leftarrow}}
\providecommand{\R}{\mathbb{R}}
\providecommand{\N}{\mathbb{N}}
\providecommand{\Z}{\mathbb{Z}}
\providecommand{\Q}{\mathbb{Q}}
\providecommand{\D}{\mathbb{D}}
\providecommand{\C}{\mathbb{C}}
\providecommand{\e}{{\mathrm{e}}}
\providecommand{\vect}[1]{\overrightarrow{#1}}
\providecommand{\vectt}[1]{\overrightarrow{\mathrm{#1}}}
\providecommand{\vvt}[1]{\overrightarrow{\mathrm{#1}}}
\providecommand{\np}[2][]{{#2~#1}}
\providecommand{\vv}[1]{{\overrightarrow{#1}}}
\renewcommand{\d}{{\mathrm{d}}}
\providecommand{\degres}{{^\circ}}
\providecommand{\degre}{{^\circ}}
\providecommand{\blue}{\color{blue}}
\providecommand{\red}{\color{red}}
\providecommand{\enskip}{\quad}
\providecommand{\llbracket}{\lbrack\!\lbrack}
\providecommand{\rrbracket}{\rbrack\!\rbrack}
\providecommand{\euro}{{\text{€}}}
\providecommand{\no}{{\text{n°}}}
\providecommand{\interval}[2]{{[#1;#2]}}
\providecommand{\intervFO}[2]{{[#1;#2[}}
\providecommand{\dotfill}{\ldots}

\providecommand{\ds}{\displaystyle}
\providecommand{\textding}[1]{#1}
\providecommand{\Oij}{\(\mathrm{O};\vec{\imath},\vec{\jmath}\)}
\providecommand{\Oijk}{\(\mathrm{O};\vec{\imath},\vec{\jmath},\vec{k}\)}
\providecommand{\Ouv}{\(\mathrm{O};\vec{u},\vec{v}\)}
\providecommand{\intervFF}[2]{\ensuremath{[#1;#2]}}
\renewcommand{\euro}{\ifmmode\text{\texteuro}\else\texteuro\fi}
\DeclareUnicodeCharacter{00B0}{\ifmmode^{\circ}\else\textdegree\fi}
\DeclareUnicodeCharacter{20AC}{\ifmmode\text{\texteuro}\else\texteuro\fi}
% Énoncés : Éducation nationale. Transcription et corrigés : APMEP, Denis Vergès et contributeurs. Découpage et classement : CoopMaths, Sébastien Lozano, Éric Elter et contributeurs. Adaptation de la présentation : Maths Explorer.
% CC BY-SA : https://coopmaths.fr/www/about/
% https://coopmaths.fr/alea/static/bac/2025/tex/bac_2025_11_sujet1_caledonie_3.tex
\begin{document}
\noindent\textit{APMEP / CoopMaths. CC BY-SA. Présentation adaptée par Maths Explorer.}\par\medskip
On considère $n$ un entier naturel non nul.

On considère la fonction $f_n$ définie sur l’intervalle $[0~;~1]$ par :

\[f_n(x) = x^n \e^{1 - x}.\]

On admet que la fonction $f_n$ est dérivable sur $[0~;~1]$ et on note $f'_n$ sa fonction dérivée.

\subsection*{Partie A}

Dans cette partie on étudie le cas où $n = 1$.

On étudie donc la fonction $f_1$ définie sur $[0~;~1]$ par :

\[f_1(x) = x\e^{1 - x}.\]

\begin{enumerate}
\item Montrer que $f'_1(x)$ est strictement positive pour tout réel $x$ de $[0~;~1[$.
\item En déduire le tableau de variations de la fonction $f_1$ sur l'intervalle $[0~;~1]$.
\item En déduire que l'équation $f_1(x)= 0,1$ admet une unique solution dans l'intervalle $[0~;~1]$.
\end{enumerate}

\subsection*{Partie B}

On considère la suite $(u_n)$ définie pour tout entier naturel $n$ non nul par 

\begin{center}$u_n = \displaystyle\int_0^1 f_n(x) \:\text{d}x$ \quad  c'est-à-dire \quad $u_n = \displaystyle\int_0^1 x^n \e^{1-x} \: \text{d}x$.\end{center}

On admet que $u_1 = \e - 2$.

\medskip

\begin{enumerate}
\item 
	\begin{enumerate}
		\item Justifier que pour tout $x \in [0~;~1]$ et pour tout entier naturel $n$ non nul,
		
		\[0 \leqslant x^{n+1} \leqslant x^n\].
		\item En déduire que pour tout entier naturel $n$ non nul,
		
		\[0 \leqslant u_{n+1} \leqslant u_n.\]
			
		\item Montrer que la suite $(u_n)$ est convergente.
	\end{enumerate}
\item 
	\begin{enumerate}
		\item À l’aide d’une intégration par parties, démontrer que pour tout entier naturel $n$ non nul on a :

\[u_{n+1} = (n + 1)u_n - 1.\]

\item On considère le script Python ci-dessous définissant la fonction \texttt{suite()} :

\begin{center}
\begin{ttfamily}
\begin{tabular}{|l|}\hline
from math import exp\\
~\\
def suite():\\
\quad u = \ldots\\
\quad for n in range (1, \ldots):\\
\quad \quad \quad u = \ldots\\
\quad return\\ \hline
\end{tabular}
\end{ttfamily}
\end{center}

Recopier et compléter le script Python ci-dessus pour que la fonction \texttt{suite()} renvoie la valeur de $\displaystyle\int_0^1 x^8 \e^{1-x} \, \text{d}x$.
	\end{enumerate}

\item 
	\begin{enumerate}
		\item Démontrer que pour tout entier naturel $n$ non nul on a :

\[u_{n} \leqslant \frac{\e}{n+1}.\]

		\item En déduire la limite de la suite $(u_n)$.
	\end{enumerate}
\end{enumerate}

\end{document}
